\chapter*{How to Use This Book}
\addcontentsline{toc}{chapter}{How to Use This Book}
\markboth{How to Use This Book}{How to Use This Book}

\textbf{Reading tiers.} Chapters are grouped into five parts. Parts~I and~II
correspond roughly to the depth expected of an entry-level (``foundation''
or ``novice''-class) license worldwide. Parts~III and part of~IV add the
depth typically expected of an intermediate (``technician'' or
``general''-class) license. The remainder of Part~IV and all of Part~V,
together with the full formula appendix, correspond to the depth expected
of the highest license class most administrations offer (often called
``advanced,'' ``extra,'' or ``full''). You do not need to read the parts
in lockstep with a specific national syllabus --- read as far as your
target license requires, then treat the rest as a reference to return to.
Part~VI is different in kind from the rest of the book: it is not exam
material at all, but an enrichment chapter applying the same radio
science to aviation, maritime, and military signals now commonly
explored with inexpensive SDR receivers.

\medskip
\noindent Throughout the book you will find four recurring boxes:

\begin{keyidea}
A concept worth memorizing verbatim, because it recurs across many later
chapters or is a favorite target of exam questions.
\end{keyidea}

\begin{examplebox}
A fully worked numerical problem, solved step by step, in the style of a
typical exam question.
\end{examplebox}

\begin{examtip}
A practical hint about how a topic tends to be tested, or a common mistake
to avoid.
\end{examtip}

\begin{safetynote}
A warning about a genuine electrical, RF-exposure, or mechanical hazard.
Safety notes are not exam trivia --- treat them as seriously as you would
a warning label on a piece of test equipment.
\end{safetynote}

\medskip
\noindent Each chapter closes with a \textbf{Chapter Review} (a condensed
list of the ideas you should be able to restate from memory) and a set of
\textbf{Practice Problems} with the reasoning needed to solve them found in
the chapter body. Appendix~A gathers every formula in the book by topic;
Appendix~B gathers SI prefixes, unit conversions, and physical constants;
Appendix~C is a glossary. Bookmark all three --- you will use them long
after you pass an exam.
